%% GENERATED from PAPER_SUBMISSION.md (pandoc) -- hand-fixed conversion; see BUILD_NOTES.md
\section{Related Work and Attribution}\label{sec:2}

We summarize prior work that our experiments directly build on. A complete attribution table for every component used in our experiments appears at the end of this section.

\textbf{SASRec} \citep{kang2018sasrec} is the base Transformer sequential recommender we re-implement. SASRec models a user's history with a causal-masked Transformer encoder and scores the next item via dot product with a learned item-embedding table.

\textbf{BERT4Rec} \citep{sun2019bert4rec} adapts BERT-style bidirectional masked-item modeling for sequential recommendation. We implement BERT4Rec as one of our 20 ablation baselines on Beauty\_and\_\allowbreak PC.

\textbf{SBERT / MiniLM} \citep{reimers2019sbert,wang2020minilm} provide pre-trained sentence embeddings. We use \texttt{sentence-transformers/}\allowbreak\texttt{all-MiniLM-L6-v2} \citep{minilm2021card} (384-d) as the default frozen item-text encoder.

\textbf{UniSRec} \citep{hou2022unisrec} is an early, widely-adopted text-based sequential recommender that uses a pre-trained language model to embed item text and learns a parametric adaptor to project these embeddings into the sequential model's hidden space; \textbf{ZESRec} \citep{ding2021zesrec} precedes it in representing items from natural-language descriptions with a pre-trained encoder for zero-shot sequential recommendation, and \textbf{RecFormer} \citep{li2023recformer} subsequently studies language representations for sequential recommendation including low-resource and cold-start evaluation --- we make no priority claim in this line. The SASRec-with-frozen-text-encoder pattern we use is a strict simplification of UniSRec.

\textbf{BLaIR} \citep{hou2024blair} (published as \citealp{hou2026blairacl}) is a RoBERTa-base model continually pretrained on 10\% of (item-metadata, review) pairs from Amazon Reviews 2023. We use \texttt{hyp1231/blair-roberta-base} (768-d) as an alternative frozen item-text encoder. Hou et al. also publish a reference implementation, \textbf{SASRecText}, that pairs SASRec with BLaIR via a 2-layer MLP adaptor \texttt{{[}768,\ 300,\ 64{]}} with Dropout(0.2) + ReLU; we adopt this adaptor design verbatim (see Appendix A.1).

\textbf{TIGER} \citep{rajput2023tiger} reformulates sequential recommendation as autoregressive generation over discrete semantic IDs produced by a residual-quantized variational autoencoder over item-text embeddings. TIGER reports strong results on Amazon Reviews 5-core benchmarks.

\textbf{LIGER} \citep{yang2024liger} extends TIGER with a dense-retrieval refinement step and reports further gains.

\textbf{Amazon Reviews 2023} \citep{hou2024blair} is the source dataset. We use the \textbf{5-core leave-last-out} protocol (described in \S\ref{sec:3}) matching the convention used by TIGER, LIGER, and similar works, which differs from the \texttt{0core\_timestamp\_w\_his\_*} HuggingFace subset used by Hou et al.'s \texttt{seq\_rec\_results/} reference implementation.

\textbf{Evaluation practice and reproducibility in recommender systems.} Our lead contribution responds to a documented evaluation-trust problem in this field. \citet{ferraridacrema2019worrying} examined 18 recent neural recommenders and could reproduce only 7 with reasonable effort, 6 of which were often outperformed by comparably simple heuristic baselines; the extended journal analysis \citep{ferraridacrema2021troubling} corroborates the pattern in greater depth and attributes it in part to the choice and optimization of the baselines used for comparison. \citet{sun2020rigorous} respond at the community level with standardized benchmarking for reproducible evaluation and fair comparison. These works diagnose the problem and propose field-level remedies; what this paper adds is complementary and \emph{per-paper}: an executable discipline that a single submission can adopt and a reviewer can re-run --- version-controlled pre-declaration of the comparator-facing confirmation (\S\ref{sec:5.2}), a fail-closed artifact gate that recomputes every artifact-gated table cell from released artifacts at build time (published comparator values are pinned as labeled external constants; \S\ref{sec:2.2} records a retirement the gate forced), local regeneration of the comparator by its own unmodified reference implementation (environment-caveated single-run regenerations, \S\ref{sec:5.6}), and symmetric publication of failures, including a pre-declaration VOIDed by its own comparability check (Appendix A.0) with a redesigned successor that subsequently passed on fresh seeds (\texttt{PREREG\_OFFICE\_V3.md}; \S\ref{sec:5.2}), the original VOID standing unchanged. None of this apparatus strengthens any empirical claim --- every non-claim of \S\ref{sec:5}--\S\ref{sec:7} stands unchanged; it certifies that the numbers printed are the numbers the artifacts recompute, under decision rules frozen before the confirmatory runs.


\subsection{Attribution table}\label{sec:2.1}

%% GENERATED by _bestrec_run/emit_latex_tables.py -- DO NOT EDIT BY HAND
% table key: table_attribution
% provenance: md lines 70-79; family=md-only (pandoc-converted; no JSON family); check={"checked": 0, "note": "md-only table; no JSON family"}
\par\medskip\begingroup\small\setlength{\tabcolsep}{4pt}\renewcommand{\arraystretch}{1.12}
\noindent
\begin{tabularx}{\linewidth}{p{0.30\linewidth}p{0.24\linewidth}Y}
\toprule
Component & Source & Where in our work \\
\midrule
SASRec architecture & \citet{kang2018sasrec} & base \texttt{SASRecSBERT} model \\
BERT4Rec & \citet{sun2019bert4rec} & one of 20 ablation baselines \\
MiniLM \texttt{all-MiniLM-L6-v2} (384-d) & \citet{reimers2019sbert}; \citet{wang2020minilm} & default frozen item-text encoder \\
BLaIR \texttt{blair-roberta-base} (768-d) & \citet{hou2024blair} (arXiv:2403.03952) & alternative frozen item-text encoder \\
2-layer MLP adaptor \texttt{{[}768,\ 300,\ 64{]}} Dropout+ReLU & \citet{hou2024blair} (SASRecText, in \texttt{external/\allowbreak AmazonReviews2023/\allowbreak seq\_rec\_results/}) & \texttt{-\/-mlp-adaptor} flag \\
Rich-text content concatenation (title + cats + brand) & derived from \citet{hou2024blair} preprocessing & \texttt{encode\_richtext\_5core.py} \\
Amazon Reviews 2023 dataset & \citet{hou2024blair} & all benchmarks \\
TIGER / LIGER & \citet{rajput2023tiger} / \citet{yang2024liger} & published comparator numbers \\
\bottomrule
\end{tabularx}
\endgroup\par\medskip



\subsection{Our additions to the above prior art}\label{sec:2.2}

We do not claim a wholly new recommender architecture. Our architectural additions are deliberately small: a causal FIR adaptation of prior frequency-filter ideas and a soft text-prototype additive term. Our contributions are:

\begin{itemize}
\tightlist
\item
  An \textbf{open-source 5-core preprocessing pipeline} (\texttt{preprocess\_5core\_standard.py}) that matches the protocol used by TIGER and LIGER (not the 0-core protocol used by Hou et al.'s repo).
\item
  A \textbf{memory-efficient chunked-full-softmax loss} that enables proper full-catalog cross-entropy training at 207k-item catalogs on a single 16 GB GPU.
\item
  An \textbf{eval optimization} (caching \texttt{all\_item\_features()} once per \texttt{evaluate()} call) that gives \textasciitilde10\ensuremath{\times} speedup at d=64 with the MLP adaptor.
\item
  A \textbf{cross-pipeline transfer study}: empirically testing which published architectural choices help on the 5-core protocol.
\item
  \textbf{Negative-result documentation} for 18 of 20 tested variants on Beauty\_and\_\allowbreak Personal\_Care 5-core.
\item
  The \textbf{competitive Video\_Games 5-core result} on the HSTU-style pure-PyTorch stack (6-seed NDCG@10 = \textbf{0.0673 \ensuremath{\pm} 0.0003}, +17.5\% over published SASRec 0.0573), shown to be architectural (ID-only \ensuremath{\approx}0.0656), with no SOTA claim. \emph{(The earlier v1-era SASRec-SBERT number 0.05509 \ensuremath{\pm} 0.00035 has been retired from Table 1a: its raw per-seed artifacts predate the artifact manifest and cannot be recomputed, so it survives only as a RETIRED provenance row in the manifest, not as a paper claim.)}
\item
  The two added components --- the \textbf{strictly causal FIR temporal filter} (cleanly isolated from an identity control under matched initialization on three categories; the older zero-gated package remains package-attributed, \S\ref{sec:5.2}) and \textbf{TAPE} --- plus the \textbf{MI frequency-5 tail case} (\S\ref{sec:5.3}) and its \textbf{interaction-thinning titration} level-contrast result (\S\ref{sec:5.4}). The novelty boundary for each is stated explicitly in Table 0 (\S\ref{sec:2.3}).
\end{itemize}


\subsection{Claim boundary}\label{sec:2.3}

Table 0 separates prior work from the paper's actual claim boundary. Frequency
filtering (including FEARec and WPGRec) and causal convolution (including Caser and
NextItNet) predate this work
\citep{du2023fearec,liu2026wpgrec,tang2018caser,yuan2019nextitnet}, as do residual
identity initialization, HSTU, frozen text encoders, and label smoothing. Our claim is limited to the
gradient-active left-causal depthwise-FIR residual used before an HSTU-style stack under
an all-position objective, together with its matched internal evaluation. The historical
zero-gated realization is singular at initialization and remains package-attributed;
only the nonsingular canonical realization supports a FIR-specific interpretation
(\S\ref{sec:3}, \S\ref{sec:5.2}).

%% GENERATED by _bestrec_run/emit_latex_tables.py -- DO NOT EDIT BY HAND
% table key: table0_novelty
% provenance: md lines 99-108; family=md-only (pandoc-converted; no JSON family); check={"checked": 0, "note": "md-only table; no JSON family"}
\par\medskip\begingroup\footnotesize\setlength{\tabcolsep}{4pt}\renewcommand{\arraystretch}{1.12}
\noindent \textbf{Table 0: Novelty boundary --- what is reused, what is changed, and how novel each piece is.}\par\smallskip
\noindent
\begin{longtable}{p{0.10\linewidth}p{0.115\linewidth}p{0.15\linewidth}p{0.185\linewidth}p{0.165\linewidth}p{0.10\linewidth}}
\toprule
Component & Prior art & What we reuse & What we change & Evidence & Novelty strength \\
\midrule
\endhead
\bottomrule
\endfoot
HSTU base & \citet{zhai2024hstu} & pointwise \texttt{silu(QK\ensuremath{^{\top}}\ +\ rab)\ V} aggregation, per-block normalization, relative attention biases & pure-PyTorch implementation; diagnosis/removal of two spurious normalizations in our initial reimplementation; equation-level verification only, core-block parity demonstrated bitwise against the reference research implementation (\S\ref{sec:3.2}); no pinned end-to-end system reproduction (\S\ref{sec:5.6}, \S\ref{sec:6.5}) & Table 1 ablations; \S\ref{sec:3.7} & none (implementation work) \\
BLaIR / SBERT text embeddings & \citet{hou2024blair}; \citet{reimers2019sbert}; \citet{wang2020minilm} & frozen encoders; Hou et al.'s MLP adaptor verbatim & nothing (frozen, off-the-shelf) & Table 1a encoder comparison & none \\
Label smoothing & \citet{szegedy2016rethinking} & standard uniform label smoothing & applied to full-catalog chunked-softmax CE in this setting & +0.0013 single-flag; +0.0012 stacked (Table 1) & none \\
Time bias & TiSASRec; HSTU \citep{zhai2024hstu} & bucketed time-interval attention bias & integration only & +0.0027 single-flag, largest classical component (Table 1) & none \\
FMLP-Rec / BSARec frequency filters & \citet{zhou2022fmlprec}; \citet{shin2024bsarec} & the idea of filtering sequential representations; frequency-domain motivation; attention-oversmoothing framing & not inserted as-is (bidirectional in their original formulations) & cited motivation (\S\ref{sec:3.7}b) & none (prior art) \\
Causal FIR adaptation (ours) & builds directly on FMLP-Rec / BSARec; causal \emph{convolutional} sequence encoders are prior art (Caser; NextItNet) & their filtering motivation and frequency framing & leak-free, left-causal depthwise FIR realization (zero-init gated) inserted before an HSTU-style stack under the all-position next-item objective, full-catalog LLOO & VG +0.0015 (5-seed); MI +0.0025, \ensuremath{\approx}80\% of the combined lift (5-seed); pre-declared breadth: IS +0.0024, CDs +0.0057 (both 5/5 seeds, CIs exclude 0, \S\ref{sec:5.2}); robust across K\ensuremath{\in}\{4,8,16,50\} & incremental \\
TAPE (ours) & TIGER \citep{rajput2023tiger}; VQ-Rec; ProtoMF & prototype / semantic-ID shared-structure motivation; frozen text embeddings & frozen k-means soft assignment gating a zero-init learnable prototype table, additive & +0.0009 single-flag (n=1); +0.0004 in the 4-seed cross-check; sub-additive (\S\ref{sec:5.1}) & incremental \\
Tail / thinning analysis (ours) & popularity-stratified evaluation (standard); long-tail SR, e.g. MELT \citep[arXiv:\allowbreak 2304.08382]{kim2023melt}; cold-start content methods, e.g. DropoutNet \citep{volkovs2017dropoutnet}, CLCRec \citep{wei2021clcrec} & train-frequency tercile bucketing; paired same-seed contrasts & dataset-conditional text-vs-ID tail contrast plus interaction- and user-mode thinning interventions with a matched-R1 (resource-controlled) design & 5-seed MI tail win (CI excludes 0); powered VG null (TOST); cross-dataset Welch p\ensuremath{\approx}0.004; user-mode dd = +0.000326 (CI excludes 0) & moderate (empirical) \\
\end{longtable}
\endgroup\par\medskip


\textbf{Closest systems I --- adjacent lines and the narrowed boundary (metadata verified 2026-07-19).} Five adjacent lines narrow what this paper may claim, and we state the narrowed boundary explicitly. (i) \emph{Text similarity for rare/unseen items:} SimRec \citep{brody2024simrec} already targets sequential-recommendation cold-start with text-derived item similarity and reports larger benefits on sparser datasets; our tail contribution is therefore graded \textbf{incremental} --- a full-catalog, cross-category AR2023 extension with intervention studies and disclosed cohort defects (\S\ref{sec:5.3}--\S\ref{sec:5.4}) --- not a new sparse-vs-dense semantic-tail idea. (ii) \emph{Convolution/filtering with attention:} C3SASR \citep{chen2022c3sasr} combines cheap causal convolutions with self-attention, AdaMCT \citep{jiang2023adamct} mixes local convolution with Transformer attention, and learnable frequency-domain filtering with gates and side-information fusion appears in DLFS-Rec \citep{liu2023dlfsrec}, DIFF \citep{kim2025diff}, TASIF \citep{luo2026tasif}, and --- with position-specific time-variant convolutional filters --- TV-Rec \citep{shin2025tvrec}; our defensible distinction is only the \textbf{gradient-active, identity-initialized left-causal depthwise FIR residual under an HSTU-style all-position objective at full catalog} --- not causal/local convolution with attention, learnable filtering, gating, or filter-content fusion as concepts. (iii) \emph{Side-information fusion scope:} NOVA \citep{liu2021nova} and DIF-SR \citep{xie2022difsr} document that direct/early fusion can be harmful and propose decoupled designs; TedRec \citep{xu2024tedrec} fuses text and ID at sequence level in the frequency domain; and AlterRec \citep{li2025alterrec} studies failures of naive ID/text fusion --- we test \textbf{only early additive frozen-text fusion}, and every ``text helps/does not help'' statement in this paper is scoped to that construction. (iv) \emph{Cold-start distinction (stated precisely):} Mecos \citep{zheng2021mecos} meta-learns for \textbf{new items with few interactions}, and RecGPT \citep{jiang2025recgpt} is a generative foundation model evaluated in \textbf{zero-shot transfer and short-history} settings --- two different cold-start settings, both distinct from our mixed zero-/low-exposure full-catalog cohort (which itself mixes zero-train-exposure and low-frequency targets, \S\ref{sec:5.3}); we make no cold-start-method comparison. (v) \emph{Reproducibility apparatus:} executable evaluation frameworks (Elliot \citep{anelli2021elliot}; DaisyRec 2.0 \citep{sun2022daisyrec}), accountability workflows \citep{bellogin2021accountability}, and reproduction audits \citep{petrov2022bert4rec} predate this work; what this paper claims is only its exact combined enforcement mechanism --- per-cell recomputation from released artifacts, provenance tripwires, and symmetric self-VOIDing adjudication --- positioned against, not above, those frameworks. TAPE's boundary likewise narrows: LLMSQRec \citep{zhang2026llmsqrec} combines frozen LLM semantic embeddings, ID embeddings, and learned prototype/codebook quantization for sparse and long-tail settings --- TAPE's continuous additive realization differs, but the broad text-prototype-ID concept is not ours. Relatedly, the chunked-full-softmax loss is implementation engineering; memory-efficient exact cross-entropy without materializing the full logit matrix is published (Cut Cross-Entropy \citep{wijmans2025cce}), and we cite it as such (\S\ref{sec:3.5}).
\textbf{Closest systems II --- null-space, frequency, and anisotropy lines (metadata verified 2026-07-20).} Five further works tighten the frozen-text/ID and frequency boundaries: \textbf{AlphaFuse} \citep{hu2025alphafuse} learns ID embeddings in the null space of language embeddings and evaluates long-tail settings --- it is now the \textbf{closest omitted frozen-text-plus-ID comparator}, and benchmarking it under our full-catalog split (or an explicit protocol-based exclusion) is queued experiment work; \textbf{DWSRec} \citep{zhang2024dwsrec} whitens anisotropic pre-trained text embeddings (dual-view); \textbf{SIDSRec} \citep{zhou2026sidsrec} separates semantic and ID channels before fusion; \textbf{BFDRec} \citep{hu2026bfdrec} does energy-aware multi-scale frequency modeling with dynamic gating (closest frequency-domain line beside TASIF/DLFS-Rec/DIFF/TV-Rec); and \textbf{ACE} \citep{lee2026ace} shows a linear-autoencoder reshaping of anisotropic LLM embeddings can itself add capacity usefully --- so our negative map's capacity-adding conclusions remain strictly \textbf{implementation- and environment-local} (they characterize our probes in our setting, not the design space). Two protocol \textbf{controls remain open experiments (disclosed, not executed):} a sequence-splitting/target-strategy parity audit mapping our all-position objective onto the taxonomy of the SIGIR 2026 reproducibility study (arXiv:2604.05309) with demonstrated comparator target-multiplicity parity, and an item-text permutation control (attribute-shuffling style, DOI 10.1145/3805712.3809839) before attributing the +2.7\% stack lift to semantic alignment rather than added representation structure.

\textbf{Closest systems III --- LLM-semantic tail and convolutional lines (metadata verified 2026-07-20).} Four further works narrow the two axes our findings sit on. \emph{Semantic-tail line:} \textbf{LLM-ESR} \citep{liu2024llmesr} (Liu, Wu, Wang, Zhang, Tian, Zheng \& Zhao --- NeurIPS 2024 spotlight) fuses LLM semantic embeddings with collaborative signals in a dual-view design for long-tail items and adds retrieval-augmented self-distillation for long-tail users; \textbf{FAERec} \citep{wei2026faerec} (Wei, Dang, Guo, Zhao \& Sun, 2026) adaptively gates ID and LLM embeddings with dual-level (contrastive + correlation) alignment for tail items; \textbf{LLM2Rec} \citep{he2025llm2rec} (He, Liu, Zhang, Ma \& Chua --- KDD 2025) trains a two-stage LLM-derived embedding model encoding both semantics and collaborative signal. LLM-ESR and FAERec target long-tail improvement directly, and LLM2Rec targets general sequential recommendation via LLM-derived embeddings --- an active methods axis. \textbf{LLM2Rec is the closest protocol overlap}: it evaluates on Amazon Reviews 2023 with 5-core filtering, leave-one-out and full-item ranking, but with maximum history length 10 (ours: 50) and a correspondingly different item universe and splits, so its numbers are not comparable under the \S\ref{sec:2.1} executable-protocol boundary; LLM-ESR, FAERec, and ConvFormer do not use this protocol. None of this changes what we claim: our tail contribution is not a method but a \emph{measured boundary} --- the MI frequency-5-heavy pattern (cross-dataset heterogeneity not established) of a deliberately frozen additive text stack against an ID-only control under a gated protocol (\S\ref{sec:5.3}), including where text does \emph{not} help (zero-exposure targets: zero hits through rank 100 in all four TFV2 categories). Whether these adaptive-fusion designs move the same frozen-protocol boundary is an open experiment we have not run, and we say so. \emph{Convolutional-sequence line:} \textbf{ConvFormer} \citep{wang2023convformer} (Wang, Lian, Wu, Li, Fan, Xu, Li \& Xie, 2023) replaces item-to-item attention with convolution in sequential user modeling --- the nearest Transformer-replacing convolutional relative alongside FMLP-Rec, BSARec, AdaMCT, and C3SASR; our FIR claim remains exactly the narrow one already stated (a gradient-active, identity-initialized leak-free left-causal FIR residual \emph{before} an HSTU-style stack under the all-position objective, supported by matched-initialization contrasts on three categories), not a claim to convolutional sequence modeling as such.

